\section{从电源键到复位向量：一次冷启动的逐级放行}

前文已说明电源、时钟、复位与 Boot ROM 的职责。这里把它们放入一次 ATX 类台式平台的冷启动 trace。该例用于解释依赖关系，不把某一块主板的具体毫秒数当成所有平台的规范：服务器、笔记本和 SoC 会使用不同信号名与管理控制器，但都必须先建立电气条件，再释放逻辑状态，最后才允许处理器取指。

\begin{pdfdiagram}[x=1cm,y=1cm]
  \node[hw state, minimum width=2.1cm] (standby) at (1.2,4.7) {待机电源有效\\EC/SIO 可以运行};
  \node[hw control, minimum width=2.1cm] (button) at (4.0,4.7) {电源键消抖\\确认启动策略};
  \node[hw control, minimum width=2.1cm] (pson) at (6.8,4.7) {拉低 PS\_ON\#\\请求主电源};
  \node[hw state, minimum width=2.1cm] (rails) at (9.6,4.7) {主电源轨爬升\\进入容差窗口};

  \node[hw state, minimum width=2.1cm] (pwrok) at (9.6,2.35) {PWR\_OK 有效\\电源声明稳定};
  \node[hw compute, minimum width=2.1cm] (vrm) at (6.8,2.35) {板级 VRM\\核心/内存电压};
  \node[hw control, minimum width=2.1cm] (clock) at (4.0,2.35) {参考时钟与 PLL\\锁定并可监测};
  \node[hw state, minimum width=2.1cm] (reset) at (1.2,2.35) {复位树分级释放\\同步到各时钟域};

  \node[hw state, minimum width=2.1cm] (strap) at (1.2,0.0) {锁存 strap\\确定启动入口};
  \node[hw memory, minimum width=2.1cm] (vector) at (4.0,0.0) {装入复位向量\\地址译码成立};
  \node[hw memory, minimum width=2.1cm] (rom) at (6.8,0.0) {Boot ROM 返回\\第一条指令字};
  \node[hw state, minimum width=2.1cm] (retire) at (9.6,0.0) {首条指令提交\\启动执行已成立};

  \draw[hw bus] (standby) -- (button);
  \draw[hw bus] (button) -- (pson);
  \draw[hw bus] (pson) -- (rails);
  \draw[hw bus] (rails) -- (pwrok);
  \draw[hw bus] (pwrok) -- (vrm);
  \draw[hw bus] (vrm) -- (clock);
  \draw[hw bus] (clock) -- (reset);
  \draw[hw bus] (reset) -- (strap);
  \draw[hw bus] (strap) -- (vector);
  \draw[hw bus] (vector) -- (rom);
  \draw[hw bus] (rom) -- (retire);
\end{pdfdiagram}
\diagramnote{ATX 类平台的一次冷启动放行链。待机域先处理按键和策略，PS\_ON\# 只请求主电源开始工作；主电源、板级 VRM、参考时钟和各级复位依次建立条件。复位释放之后还要成功锁存启动配置、命中复位向量、从不可变入口取得合法指令并完成首条提交，才能说处理器真正开始执行。}

\topic{“有电”至少包含四个不同事实}

插座有输入、待机轨有效、主电源输出存在、负载点电压满足处理器工作窗口，是四个不同层次。待机 LED 点亮只证明前两层的一部分；风扇转动只证明某条电源路径和控制输出工作；即使万用表读到额定直流电压，也不能证明爬升斜率、纹波、短时跌落和 PGOOD 时序满足要求。启动诊断必须把这些证据分开。

待机电源让 EC、Super I/O 或管理控制器在主 CPU 断电时仍能扫描电源键、处理唤醒源并执行策略。按键通常是低速机械输入，要先消抖，还可能受长按强制关机、机箱盖、温度、适配器能力或上次故障锁存影响。控制器接受启动请求后才断言主动低的 PS\_ON\#；这不是“CPU 复位释放”信号，只是主电源开始建立输出的请求。

\Needspace{0.40\textheight}
\begin{longtable}{L{0.08\linewidth}L{0.2\linewidth}L{0.27\linewidth}L{0.33\linewidth}}
\toprule
阶段 & 主要观察点 & 放行条件 & 若停在这里通常检查 \\
\midrule
P0 & 待机轨、EC 时钟与复位 & 待机电压稳定，控制器能取指 & 输入电源、保险/保护、待机稳压 \\\tablerowrule
P1 & 按键电平与事件日志 & 消抖通过，策略允许启动 & 按键、上拉、EC 固件和故障锁存 \\\tablerowrule
P2 & PS\_ON\# 与主电源轨 & 电源接受请求，各轨进入容差 & 过流、短路、连接器和电源保护 \\\tablerowrule
P3 & PWR\_OK、板级 enable & 主电源声明稳定，二级电源可启动 & 纹波、斜率、掉电与时序窗口 \\\tablerowrule
P4 & VRM 输出与 PGOOD & 核心、内存、I/O 电压均合格 & VRM 相位、负载短路、遥测和温度 \\\tablerowrule
P5 & 晶振、参考时钟、PLL lock & 目标域有连续合规时钟 & 晶振、时钟缓冲、配置 strap \\\tablerowrule
P6 & 各复位端与同步器 & 依赖域就绪，复位按顺序同步释放 & reset supervisor、CPLD/EC 和跨域同步 \\\tablerowrule
P7 & SPI/ROM 地址与数据活动 & 复位向量命中只读启动入口 & strap、地址别名、Flash 供电和信号完整性 \\\tablerowrule
P8 & 首取指、异常与进度码 & 指令合法且首个架构状态提交 & Boot ROM 内容、总线错误和处理器故障 \\
\bottomrule
\end{longtable}

\topic{PGOOD 是声明，不是能量本身}

PWR\_OK 或板级 PGOOD 表示监测电路认为相关电压已经进入允许窗口并维持了规定条件。逻辑应把它当作放行条件之一，而不能用固定延时替代全部监测：轻载时电压可能很快到位，短路或振荡时等再久也不会安全。反过来，PGOOD 线路被错误上拉也可能在电压尚未成立时给出假许可，所以验证要同时测模拟电压和数字状态。

主电源轨稳定后，板上 VRM 才为 CPU 核心、SoC、内存和其他负载点建立更低电压。多个电源域存在依赖：I/O 电源先于核心或相反、内存参考电压跟随、PLL 模拟电源低噪声等要求由平台时序控制器落实。每个域的 PGOOD 汇聚成复位树的条件；任何一个必需域掉出窗口，都应重新断言相应复位，而不是让该域继续执行未知状态。

\topic{复位释放也是跨时钟域事务}

复位的断言通常要求快速把逻辑带回安全状态，释放则要同步到目标时钟域，避免不同触发器在不同边沿离开复位而产生亚稳态。时钟存在并不等于 PLL 已锁定，PLL lock 有效也不保证所有分频和门控已经打开。平台常先释放时钟/电源控制逻辑，再释放内存控制器、互连和处理器，最后释放依赖它们的设备。

strap 的采样窗口也位于这条边界附近。板级上拉/下拉必须在采样沿之前稳定，复位释放后锁存结果才成为本次启动的配置代际。软件运行时读到的 strap 镜像只是在解释已经发生的选择；修改相关寄存器通常不能 retroactively 改变本次复位入口。

\topic{复位向量命中仍不等于首条指令完成}

处理器退出复位后把取指地址置为架构规定的入口。地址译码必须把该地址导向片上 ROM、SPI 映射窗口或其他可信启动介质；存储端返回正确编码，取指端通过校验和权限检查，译码器接受该编码，首条指令才进入执行。若第一条就是跳转，还要看到目标地址的下一次总线活动；若首条触发异常，则“总线上出现过复位向量”不能证明启动代码正常推进。

\Needspace{0.34\textheight}
\begin{longtable}{L{0.22\linewidth}L{0.25\linewidth}L{0.25\linewidth}L{0.18\linewidth}}
\toprule
表面现象 & 已经证明 & 尚未证明 & 下一证据 \\
\midrule
待机灯亮 & 至少一部分待机供电存在 & EC 正在执行、主电源可启动 & 待机纹波、EC 时钟/按键日志 \\\tablerowrule
风扇短转后停止 & 主电源曾被请求或风扇被脉冲驱动 & 所有主轨、VRM 和复位均合格 & PS\_ON\#、PWR\_OK 与故障锁存 \\\tablerowrule
核心电压存在 & VRM 有输出 & 时钟锁定、CPU 已退出复位 & PGOOD、时钟和复位端 \\\tablerowrule
SPI 有片选与时钟 & 某启动主机正在读取 & 地址正确、数据合法、CPU 会提交 & 地址/数据解码、进度码 \\\tablerowrule
串口无输出 & 只能说明未观察到串口字符 & CPU 必然未执行 & 首取指 trace、早期 GPIO/POST 码 \\
\bottomrule
\end{longtable}

\topic{掉电和失败恢复必须回到已知代际}

某级超时后，控制器应记录最后成功阶段和首个失败原因，按平台允许的顺序重新断言复位、关闭主电源并等待储能释放，再开始新的启动代际。只把 CPU reset 拉一下，不能清除卡死的 VRM、时钟发生器、互连或外设状态；频繁自动重试还可能掩盖过流和热故障。新的代际要让旧的 PGOOD、完成脉冲和启动定时器失效，避免迟到信号误放行。

示波器、逻辑分析仪、EC/BMC 事件日志、主板 POST 码和处理器 trace 各覆盖不同层次。最有效的排查不是随机替换部件，而是沿 P0 到 P8 找到“最后一个可靠事实”和“第一个未成立条件”。这与全书反复使用的事务方法一致：请求已经发出，不代表对方接受；电压已经出现，不代表复位可以释放；复位已经释放，也不代表首条指令已经提交。
